% !TeX root = ../../Thesis.tex
\section{Computational Methodology}\label{chap2}
This chapter documents the numerical framework in STAR-CCM+, 
detailing the one-fluid VOF method for separated free surfaces and 
the Eulerian Multiphase model with Large-Scale Interface capabilities 
for regime-spanning applications. It explains interface-capturing, 2D adaptive remeshing to sharpen fronts, 
and the closure palette for surface tension using the continuum surface theory, 
interfacial momentum (drag, lift, turbulent dispersion, wall lubrication), 
and interphase heat transfer. Finally, it states modeling assumptions and boundary conditions for both 
infiltration and the envisioned flashing application.
\subsection{Multiphase-Flows in STAR-CCM+}
STAR-CCM+ employs a finite-volume, cell-centered scheme with 
 variables stored at cell centroids and interpolated to faces, 
 using second-order upwind for convection, least-squares gradients 
 with Venkatakrishnan limiting, and the Semi-Implicit Method for Pressure Linked Equations 
 (SIMPLE) scheme for segregated pressure-velocity coupling. This ensures a balance of accuracy, boundedness, and 
 robustness in the presence of large density jumps and under-resolved interfaces. 
 Interface-capturing methods necessarily produce a minimum interface thickness of a few cells, motivating bounded convective schemes 
 and careful limiter usage to control numerical diffusion and spurious currents. RANS modeling is used with the SST $k-\omega$ model, realizability controls, turbulence intensity and length-scale boundary specification, and optional quadratic/cubic 
stress-strain relations for anisotropy in flows with strong secondary motions

At large scale interfaces, the near-interface flow often exhibits
 laminar characteristics with the development of boundary layers. However, in situations
 with significant velocity gradients, turbulent fluctuations can arise in the
 regions adjacent to the interface (especially when the interface is mobile or
 highly distorted). Recognizing that standard RANS turbulence
 models tend to overpredict turbulence intensity near phase boundaries,
 STAR-CCM+ allows the user to apply turbulence damping strategies
 (vorticity limiters) which suppress excessive eddy viscosity and
 restore physically consistent laminar sublayers adjacent to the interface.

\subsubsection{Volume-of-Fluid Method}
The volume-of-fluid (VOF) method is an interface capturing method that utilises the governing 
 equations of the one-fluid model as detailed in~\sect{sec:one-fluid}. The method presumes that
 the computational mesh is fine enough to resolve the interface topology.
 Using the volume fraction as the marker function - the interface is assumed to be present
 in cells with volume fraction values between 0 and 1.
 The distribution of the phase \(\ind{i}\) is driven by the volume fraction transport equation,
\begin{equation}\label{eq_vof_transport}
\pderiv{}{t}\integ{V}{} \alpha_{\ind{i}}\,\diff{V} +
\ointeg{\partial V}{} \alpha_{\ind{i}} \vect{v}_{\ind{i}} \cdot \diff{\vect{a}} = 
\integ{V}{} \left(S^{\alpha}_{\ind{i}} - \frac{\alpha_{\ind{i}}}{\rho_{\ind{i}}} \frac{D \rho_{\ind{i}}}{D t}\right)\,\diff{V}
- \integ{V}{} \frac{1}{\rho_{\ind{i}}} \nabla \cdot \left(\alpha_{\ind{i}} \rho_{\ind{i}} \vect{v}_{d,\ind{i}}\right)\diff{V}\,.
\end{equation}
 In case of two phases, the volume fraction transport is solved only for the first phase, thus in each volume
 element, the volume fraction of phase \(j\) is given as \(\alpha_{j} = 1 - \alpha_{\ind{i}}\).
 
For convection of volume fraction, STAR-CCM+ employs a modified High-Resolution Interface 
Capturing (HRIC) scheme that blends downwind and
 upwind via normalized-variable 
diagrams with Courant-number corrections to preserve boundedness and sharpness.The spurious oscillations that may arise in the volume fraction field due to meshes with large aspect ratios
 are suppressed by artificially smoothening the volume fraction gradient across the interface.~\cite{star-ccm+}

\subsubsection{Eulerian Multiphase Method}
Star-CCM+ implements the governing equations of the two-fluid model as detailed in~\sect{sec:two-fluid} as the eulerian multiphase (EMP) model. 
 The two-fluid model is predominantly used to simulate dispersed flows and the mesh resolution is
 assumed to be coarse enough to not capture the interface topology between the phases. The EMP model,
 however, utilises the multiple flow regime framework to allow the simulation of multiscale flows.
 Based on increasing volume fraction of the
 secondary phase, \(\alpha_{\ind{s}}\), in a finite volume cell,
 the topology transitions through three principal 
 regimes~\cite{gadaLargeScaleInterface2017, star-ccm+}, where the secondary phase is:
 \begin{itemize}
    \item dispersed in the primary phase (first dispersed regime, \(0 \le \alpha_{\ind{s}} \le \alpha_{\ind{fr}} \))
    \item separated from the primary phase through a large scale interface (interfacial flow regime, \(\alpha_{\ind{fr}} \le \alpha_{\ind{s}} \le \alpha_{\ind{sr}} \)) 
    \item continuous while having the primary phase dispersed (second dispersed regime, \(\alpha_{\ind{sr}} \le \alpha_{\ind{s}} \le 1\))
 \end{itemize}
The first and second regime terminuses, \(\alpha_{\ind{fr}}\) and \(\alpha_{\ind{sr}}\), mark the values of \(\alpha_{\ind{s}}\)
that define the end of the first dispersed and large scale interface regimes respectively. STAR-CCM+ uses regime-specific closure models for 
 the interfacial momentum and energy terms. These models are blended using a weighting
 scheme according to the calculated proportion of each regime in the cell.

Adaptive Interface Sharpening (ADIS) is used 
in the multiregime framework to switch to HRIC near large interfaces and 
revert to upwind away from interfaces, avoiding spurious
 sharpening in smooth regions while limiting oscillations at high aspect-ratio meshes.

\subsubsection{Adaptive Mesh Refinement}
STAR-CCM+ provides provisions to include adaptive mesh refinement to locally refine the mesh 
 based on either pre-defined or user-field functions. For 2D meshes, this provision has
 not been included, however the Remeshing solver can be included in the solution to
 locally refine the mesh during the evolution of the solution. The mesh refinement 
 is unidirectional and does not coarsen already refined regions of the mesh.

\begin{figure}
  \centering
  \includegraphics[width = 0.9\linewidth]{Chapters/Chap2/figs/sub1_amr} 
\caption{Adaptive mesh refinement based on volume fraction of liquid as refinement criteria.} \label{fig:amr}
\end{figure}
In context of this thesis, a user defined field function based on the volume fraction of water 
has been used in the implementation of the VOF method to refine the mesh 
with increasing time steps.~\fig{fig:amr} displays progressive refinements in the remeshing 
process as the solution evolves for the case of water flowing into a chamber. 
%"neither of the two methodology - VOF or two-fluid - is suitable for multiscale flows~\cite{gadaLargeScaleInterface2017}".
%"most data for the semi-empirical relations covers rather narrow domains~\cite{yadigarogluNatureMultiphaseFlows2018}"
%"if one considers a non flow situation in the two horizontal layers, the forces acting on the 
%two fluids and their interface are gravity and surface tension - in this case we speak of
%the RT instability~\cite{hewittModellingStrategiesTwoPhase2018}."
%"there exists a wide spectrum of scales, ranging from the global one dictated by the flow
%configuration, to the smallest one where turbulent energy is dissipated into heat\cite{shyyComputationalTechniquesComplex1997}"

\subsection{Closures}\label{sec:closures}
Closures are critical to the accurate prediction of the relevant flow variables and the distribution
 of volume fraction~\cite{sugrueRobustMomentumClosure}. Existing literature reviews 
 the overwhelming number of closures available that yield a range of varying performance 
 along with their comparison with experimental data, however literature
 presenting CFD simulations tend to choose 'result-oriented' closures that aren't
 transferable to a wider range of flow situations.
 Literature admits that the industry relies on the tunable parameters in the closures to 
 match the specific experimental data being simulated or the arbitrariness of this choice
 to get improved flow predictions.~\cite{krepperValidationClosureModel2018, sugrueAssessmentSimplifiedSet2017, liaoEulerianModellingTurbulent2018}
 
Models may also provide satisfying predictions for certain parameters or situations while
 failing in others and a general inconsenus prevails while discussing the effects of interfacial
 forces. Early work has tended to neglect 
 the effects of lift or virtual mass forces, or chosen to use constant coefficients to 
 describe them while using formulations for the drag forces based on 
 relevant non-dimensional numbers.
 The effect of non-drag forces, however, tends to be as significant on the 
 distribution of phases.~\cite{liaoEulerianModellingTurbulent2018}

 Since it is 
 difficult to separate the effects of the various forces in actual experiments,
 the development of these closures assumes certain idealisations when they're not
 partly or wholly based on empirical data, rendering them highly flow 
 dependent. The methodology of resolving a 
 complex interfacial momentum transfer through a linear superposition 
 of several simple forces is also questioned.~\cite{sugrueRobustMomentumClosure}. 

 Closures should reflect the local physical phenomenon occurring on the non-resolved scale. 
 But this requires a better understanding of the physics of such local flow phenonema. 
 Consequently, the closures should reflect the physics including their dependence 
 on material parameters and local flow characteristics. Best practices suggest that all relevant closure models should be included 
 - if the model reflects the physics accurately, 
 the locally unimportant closures will be blended out automatically.

The industry has, however, attempted to establish generality that can be transferred 
 to multiple flow situations and has proposed baseline frameworks. There is a general consensus
 that the role of turbulence is underestimated~\cite{colomboBenchmarkingCFDModelling};
 current models also aim to reduce 
 uncertainty and inconsistency in model selection~\cite{colomboBenchmarkingCFDModelling}, 
 adopt closures that avoid "case-specific adjustments to unphysical coefficients and 
 tunable parameters"~\cite{sugrueRobustMomentumClosure}, and utilise more
 physics based closures that are more physically consistent~\cite{liaoEulerianModellingTurbulent2018,colomboBenchmarkingCFDModelling}.

 \subsubsection{Surface Tension}
Surface tension, although an interfacial force, is implemented using the 
 Continuum Surface Force approach as a volumetric force in the momentum equation.~\cite{star-ccm+}
 Divided into normal and tangential components and using the mean 
 curvature of the free surface 
 \(\kappa = - \nabla \cdot \left(\nabla \alpha_{\ind{i}}/|\nabla \alpha_{\ind{i}}|\right)\),
 the volumetric surface tension force is given as 
 \begin{equation}
   \vect{f}_{\sigma} = \underbrace{-\sigma \nabla \cdot \frac{\nabla \alpha_{\ind{i}}}{|\nabla \alpha_{\ind{i}}|} \nabla \alpha_{\ind{i}}}_{\vect{f}_{\sigma,n}}
    + \underbrace{(\nabla \sigma)_{t}|\nabla \alpha_{\ind{i}}|}_{\vect{f}_{\sigma,t}}
 \end{equation}
 The EMP treats surface tension by splitting it among the phases occupying
 the cell using the local volume fractions. The pressure gradient
 is then the summation of the 
 individual phase surface tension contributions \(\vect{f}_{\sigma, \ind{i}}\)~\cite{gadaLargeScaleInterface2017,star-ccm+},
 \begin{equation}
   \vect{f}_{\sigma, \ind{i}} = \alpha_{\ind{i}} \sigma \kappa \nabla \alpha_{\ind{i}}.
 \end{equation}

\subsubsection{Interphase Momentum Closures}
Accurate interfacial momentum exchange controls phase distribution and slip, so 
 closure choices strongly influence mean velocities and void-fraction fields.

Drag force is generally more dominant than other forces as they tend to be a magnitude
 higher than other forces~\cite{tas-koehlerCriticalAnalysisDrag2021a}. Most drag models perform 
 equally well at low gas volume fractions. However as they are based on ideal conditions, 
 they do not account for bubble swarm effects at higher gas volume fractions and show reduced 
 accuracy. Certain swarm corrections do exist, but they are averaged along a radial profile when
 the swarm effects are also felt in the bulk~\cite{sugrueRobustMomentumClosure, sugrueAssessmentSimplifiedSet2017,colomboBenchmarkingCFDModelling}.
 The accurate modelling of interaction area density and length becomes relevant for the calculation
 of the drag force. 
 Using the linearised drag coefficient, \(A_{D}\), and the relative velocity between the phases
 \(\vect{v}_{r} = \vect{v}_{j} - \vect{v}_{i}\), STAR-CCM+ implements the drag force as,
 \begin{equation}
   \vect{F}_{ij}^{D} = A_{D}\vect{v}_{r}
 \end{equation}
 The linearised drag coefficient is related to the standard definitions of the drag 
 coefficient \(C_{D}\) using
 \begin{equation}
   A_{D} = C_{D} \frac{1}{2} \rho_{c} |\vect{v}_{r}| \frac{a_{cd}}{4}
 \end{equation}
 where \(a_{cd}\) is the interfacial area density and \(a_{cd}/4\) represents the projected area of
 the equivalent spherical particle. 

The lift force models are particularly sensitive and uncertain in their performance~\cite{sugrueAssessmentSimplifiedSet2017}.
 They suffer from the same discrepancy between performance in low and 
 high gas volume fractions. Moreover, existing 
 lift force implementations in STAR-CCM+ fail to replicate the experimental data
 they were derived from~\cite{sugrueRobustMomentumClosure}.
 Assuming \(\alpha_{p}\) is the volume fraction of the primary (continuous)
 phase, STAR-CCM+ implements the lift force as
 \begin{equation}
   F_{ij}^{L} = C_{L,\text{eff}} \alpha_{d} \rho_{p} 
   \left[\vect{v}_{r} \times \left(\nabla \times \vect{v}_{p}\right)\right].
 \end{equation} 

A turbulent dispersion force is supposed to act on the gas phase and 
 spreads this dispersed phase due to turbulence
 in the continuous phase. This is typically modelled as 
 as a diffusion term proportional to the turbulent 
 kinetic energy gradients~\cite{liaoFlashingEvaporationDifferent2013}. 
 The force is theorised to arise from the averaging of the fluctuations
 of the bubbles due to the turbulence.
 \begin{equation}
   F_{ij}^{TD} = A_{D} \vect{v}_{TD}
   = A_{D}\,\underbrace{C_{0} \frac{\nu_{c}^{t}}{\sigma_{\alpha}}\tens{I}}_{\tens{D}_{TD}} \cdot 
   \left[
       \nabla \ln(\alpha_{d}) - \nabla \ln(\alpha_{c})
       \right]
 \end{equation}
 The closure relation by Burns~\cite{star-ccm+} uses the Favre averaging of the drag
 force and can be expressed as
 \begin{equation}
   F_{ij}^{TD} = C_{D} C_{TD} \frac{\mu_{l}^{t}}{\rho_{c} \sigma_{TD}}\vect{v}_{TD},
 \end{equation}
 where \(C_{D}\) is the drag coefficient, \(C_{TD}\) is an adjustable turbulent dispersion
 calibration coefficient, \(\mu_{l}^{t}\) is the liquid turbulent viscosity, and
 \(\sigma_{TD}\) is the turbulent Prandtl number.

Most lift
 coefficient correlations don't account for the decrease of the magnitude of the
 lift force to zero 
 near the walls. This low gas volume fraction at the walls has historically
 been explained
 through a liquid film forming at the walls repelling the bubbles. This has been demonstrated to
 not actually occur in practice. The wall lubrication force has thus been criticised
 as lacking any physical justification~\cite{sugrueRobustMomentumClosure,colomboBenchmarkingCFDModelling}.
 Newer closures~\cite{sugrueRobustMomentumClosure,star-ccm+} utilise a turbulent dispersion force based 
 formualtion to maintain the desired gas volume fraction and use the closure instead
 as a correction term to resolve the phase distribution at the boundary lost due to
 the averaging process. STAR-CCM+'s implementation of the wall lubrication force
 \begin{equation}
   F_{ij}^{WL} = - C_{WL}(y_{w})\alpha_{s} \rho_{p} 
   \frac{|\vect{v}_{r} - (\vect{v}\cdot \vect{n})\vect{n}|}{d_{b}} \vect{n}
 \end{equation}
 can be changed to implement alternative models instead by defining user defined 
 field functions for \(C_{WL}\).

\subsubsection{Interphase Heat Transfers}
STAR-CCM+ models heat transfer by using the analogy to Newton's law of cooling. This entails the 
estimations of accurate heat transfer coefficient correlations that can link the 
liquid superheat to vapour generation in a reliable manner. This choice controls the predictive
accuracy of the rate of phase change under non-equilibrium conditions. As with the case of
momentum closures, there is no universally valid correlation as the law of cooling is a simplification
of the actual heat transfer mechanisms in multiphase flows. A composite heat transfer coefficient
correlation that accounts for the conduction- and convection-based limits can be constructed using 
a linear superposition of the various individual heat transfer correlations. Heat transfer through
turbulence can also be similarly considered. 

\subsubsection{Nucleation and Interphase Transfer Closures}
Homogeneous nucleation and subsequent bubble growth are not inherent to the two-fluid model's 
fundamental equations. To capture these effects, a scalar transport equation for the bubble number density must be solved along with the flow variables, with source terms accounting for:
\begin{enumerate}
    \item Nucleation: A nucleation model, such as the one based on the classical nucleation theory (CNT), provides the source of new bubbles. This term depends exponentially on the liquid superheat.
    \item Bubble Growth/Collapse: The growth of these newly formed bubbles is governed by interphase heat and mass transfer. The mass transfer rate is given by a model like the 
    HKS equation, while the heat transfer rate requires a closure for the interphase heat transfer coefficient.
    \item Coalescence and Breakup: To account for changes in bubble size distribution, models for bubble coalescence and breakup are also required. 
    These models are typically empirical and depend on the local turbulence and velocity gradients. In the context of our experimental conditions,
    as bubbles in low gravity show lower detachment rates and grow much larger in size than in terrestrial conditions. Coalescence will
    dominate as growing bubbles eat each other up but do not leave the nucleating surface. This factor, due to 
    the modelling complications can be ignored. 
\end{enumerate}
While a simple, pressure-driven mass transfer model like Lee model can be used to bypass these complexities, its parameters are often highly tuned and lack physical basis. 
The population balance approach, while more complex, provides a more fundamental and physically consistent framework for modeling two-phase flows with a changing bubble size distribution.

\subsection{Model of the Experiment}
The experimental flow system encompasses distinct flow regimes that 
differ in physical nature and complexity across its various components. 
To effectively capture these differences, the numerical simulation is 
logically divided into two subproblems: the liquid intrusion phase, 
characterized by transient multiphase displacement of air by water, and 
the subsequent evaporation phase, involving rapid phase change under vacuum 
conditions.

In order to establish kinematic and dynamic similarity with the experimental 
setup, the modelled domain is geometrically same as the actual setup.
Relevant scales are defined and STAR-CCM+'s automatic assignment
of a 1m depth to 2D domains is accounted for in the initial conditions 
for the infiltration problem and establish dynamical similarity. 

It is acknowledged that the current modeling employs two-dimensional 
simulations predominantly for computational efficiency, which inherently 
introduces limitations in capturing fully three-dimensional flow structures 
and turbulence anisotropies present in the actual experimental setup. 
Similarly, the adoption of the Boussinesq approximation within the SST 
$k$-$\omega$ turbulence model introduces simplifications that neglect density 
variations except in buoyancy terms, potentially diminishing accuracy 
in regimes where density gradients or compressibility effects become significant.

Despite these simplifications and inherent approximations, the staged 
subproblem approach facilitates focused investigation of the dominant 
physical phenomena in each phase. The following subsection details 
the modeling setup specific to the liquid intrusion phase, outlining 
assumptions, governing equations, boundary conditions, and validation strategies.

\subsubsection{Assumptions}
We model the experiment on a 2D domain. Water is assumed to be 
incompressible in both the subproblems. Air is modelled as an 
ideal gas in the infiltration problem. All phases are assumed to be immiscible.
We assume that the continuum hypothesis holds and that the macroscopic 
averaged quantities are smooth. 
We assume transport properties such as viscosity and conductivity are constant 
over the averaging procedure and that interfaces are infinitely 
thin. 

It is assumed that the fundamental assumptions in eddy-viscosity modelling are
true and that Reynolds stress tensor is proportional 
to the mean velocity strain rate through an eddy viscosity defined as in~\eq{eq:turb_visc}~\cite{star-ccm+}.
\begin{equation}
  \tens{T}_{\text{RANS}} = \mu_{t} (\nabla \vect{v} + \nabla \vect{v}^{T}) - \frac{2}{3}(\mu_{t} \nabla \vect{v}) \tens{I}\,,
\end{equation}
In cases of continuous-dispersed regimes, we
assume that the interfacial shear stress is equivalent to the shear stress in the continuous phase.

In the formulation of the liquid infiltration problem, no relative velocity is considered 
between the phases and the diffusion velocity in~\eq{eq:v-diff} is 0. Apart from the 
liquid inflow, no additional source terms are present. We assume that the heat exchanged 
between the environment and the boundaries of the system is dominated by conduction
and prescribe a very low heat transfer coefficient to account for the heat lost to the environment.
This prescription is not considered in the evaporation problem, where the system is assumed 
to be isolated.

\subsubsection{Governing Equations and Boundary Conditions}
The infiltration problem uses the governing equations for the one-fluid formulation
 and the evaporation equation utilises the two-fluid formulation
 as described in~\sect{sec:one-fluid} and~\sect{sec:two-fluid} respectively.
 As only the 2D aspects of the problem are considered, 
 the momentum and energy equations in the $z-$axis directions
 are not solved.

\subsubsection*{Boundary Conditions: Liquid Infiltration}\par\vspace{-\baselineskip}\vspace{\parskip}
The jump conditions at the interface without phase change are considered based on the discussion~\sect{sec:jumps}.
These are summarised below:
\begin{subequations}
  \begin{align}
    \left[\! \left[\vect{v} \right]\! \right] &= 0\,, \\
    \left[\! \left[ \left(\mu \left(\nabla \vect{v} + (\nabla \vect{v})^{T} \right) - (p + \rho \vect{g} \cdot \vect{x}) \tens{I} 
    \right) \cdot \vect{n} \right] \! \right] &= \sigma \kappa \vect{n} + \left[\! \left[\rho \right]\! \right] \vect{g}\,.
  \end{align}
\end{subequations}
The domain boundaries of the system are specified as walls without slip, thus 
 fluid velocity at wall is 0. The inlet is assigned as a velocity inlet, with
 the velocity and temperature of the inflowing water specified.

The boundary specifications for heat transfer for the infiltration problem are summarised below:
 \begin{subequations}
  \begin{align}
    &\text{Walls: Convection specification, } && \dot{q} = h_{\text{ext}} (T_{w} - T_{\infty})\,,\\
    &\text{Outlet: Adiabatic} && \dot{q} = 0\,,\\
    &\text{Inlet: Specified temperature of infiltrating water} && T = T_{in}\,.
  \end{align}
\end{subequations}
For the specification of the wall temperature $T_{\ind{w}}$ is the static fluid temperature at the wall.

\subsection{Summary}
The chosen discretisation, advection controls, and 
remeshing strategy balance boundedness, interface sharpness, 
and stability in the presence of large density jumps and moving free surfaces.
 While the approach is assumed to be pragmatic, 
 closure selections and additional damping options aim to mitigate known artifacts near phase boundaries, 
 enabling credible verification tests and 
a setup applicable for the simulation of the infiltration problem.